  

%%%%%%%%%%%%%%Empieza la presentacion%%%%%%%%%%%%%%%%%%%%%%%%%%%%%%%%%%%%

%\begin{frame}
%\frametitle{Esquema}
 %   \tableofcontents
%\end{frame}

\section{Objetivos}

\begin{frame}
\frametitle{Objetivos del tema}

Caracterizar matemáticamente:

\begin{itemize}

\visible<2->{
\item El muestreo y reconstrucción en el dominio de la frecuencia, a través de la Transformada de Fourier Discreta (DFT)
}

\visible<3->{
\item La relación entre la DFT y las series y transformadas vistas en temas anteriores (transformada Z, Serie de Fourier, Transformada de Fourier).}

\visible<4->{
\item Las propiedades de la DFT.}




\visible<5->{\item Comportamiento de los sistemas LTI bajo la DFT.
}

\end{itemize}

\visible<6->{
Introducir los métodos de computo eficiente asociadas a la DFT:

\begin{itemize}

\visible<7->{
\item Transformada Rápida de Fourier.}

\visible<8->{
\item Respuesta de los sistemas LTI.}
\end{itemize}
}


\end{frame}

